% results.tex (v10). Numbers are macros (shared_v10). Delta = RMSE(first) - RMSE(second), 95% block-bootstrap interval in brackets; negative favours the first method.
\section{Results}\label{sec:results}

\subsection{Observations and the year-specific climate input}
The label set contains \nncyAk{} cell-years in Alaska (\nncellAk{} cells, \nyrsAk), \nncyLe{} in the Lena Delta (\nncellLe{} cells, \nyrsLe) and \nncyCa{} in north-western Canada (\nncellCa{} cells, \nyrsCa) (Fig.~1a,b and Table~1). The years are unevenly covered: \ntopyrshareAk{}\% of the Alaska cell-years were observed in \ntopyrAk, whereas the Lena Delta labels are spread over \nnyrLe{} years. Most observations were made in August (median day of year \ndoymedAk{} in Alaska, \ndoymedLe{} in the Lena Delta and \ndoymedCa{} in Canada), before the end of the thaw season, and thaw continues after the observation day (Supplementary Fig.~S1). The thawing degree-days accumulated to the observation day were therefore smaller than the full-season sum, with median ratios of \ntddfracAk, \ntddfracLe{} and \ntddfracCa{} (Fig.~1d). All models in this study use the climate of the observation year from ERA5-Land monthly data for 2010--2024, accumulated to the observation day; the comparison input is the 2011--2024 mean of the same variables, which corresponds to the multi-year inputs of existing products (Fig.~1c,d). Fig.~2 summarises the physics-residual workflow, the four model structures compared and the five validation designs.

Replacing the mean climate input by the year-specific input reduced the spatial-block error in the Lena Delta for the recalibrated Stefan model (\ndforcLePo{} cm) and for the physics-residual model (\ndforcLeRo{} cm), and in Canada for the physics-residual model (\ndforcCaRo{} cm) (Fig.~3a). In Alaska the change was not resolved (Stefan \ndforcAkPo{} cm), because most labels come from a single year and the year-specific input then adds little information beyond the spatial pattern. The year-specific input is also what makes year-labelled maps possible: every prediction refers to the climate of a stated year.

\subsection{Validation design and transfer to a withheld region}
The ranking of methods depended on the validation design (Fig.~3b). Random cross-validation by cell gave the lowest errors for direct ML (Alaska \nrmseAkrandDz{} cm), and spatial-block cross-validation with 0.5$^\circ$ blocks raised the direct-ML error by \noptimAkDz, \noptimLeDz{} and \noptimCaDz{} cm in Alaska, the Lena Delta and Canada. Holding out whole years raised the direct-ML error in the Lena Delta from \nrmseLeblkDz{} to \nrmseLeyrDz{} cm, because its labels are spread over many years, and holding out blocks and years together gave the highest within-region errors of all methods in all three regions. When a whole region was withheld and the models were trained on the other two regions, the direct-ML error in the Lena Delta rose to \nrmseLeregDz{} cm, compared with \nrmseLeregPz{} cm for the Stefan model with the coefficient of the other regions (difference \ndrefLeregDz{} cm). The physics-residual model stayed close to Stefan in the same test (\ndrefLeregRo{} cm) and was better than Stefan in Canada (\ndrefCaregRo{} cm). Spatial-block cross-validation is therefore used as the main design below, and region hold-out as the test of transfer.

\subsection{Gain of the physics-residual model and its conditions}
Under spatial-block cross-validation with the year-specific input (2010--2024 labels), the physics-residual model reduced the error of the recalibrated Stefan model from \nrmseLeblkPo{} to \nrmseLeblkRo{} cm in the Lena Delta (\ndrefLeblkRo{} cm) and from \nrmseCablkPo{} to \nrmseCablkRo{} cm in Canada (\ndrefCablkRo{} cm). In Alaska both models reached about \nrmseAkblkRo{} cm (\ndrefAkblkRo{} cm). The gain required terrain and soil covariates (Fig.~3c). With climate covariates only, the residual model did not improve on Stefan in Alaska (\ndfeatAkFclim{} cm) and Canada (\ndfeatCaFclim{} cm); adding terrain and soil gave the gains above, whereas adding MODIS vegetation and snow timing (2010--2024) or high-resolution soil and land-cover classes changed the result by less than 0.5 cm in the Lena Delta and Canada; in Alaska the high-resolution soil and land-cover classes added a small gain (\ndfeatAkFhires{} cm).

The position of the physical information mattered (Supplementary Fig.~S2). Adding pseudo-labels from the Stefan model to the training set reproduced the Stefan model (Lena Delta \ndrefLeblkDo{} cm), and using the Stefan prediction as an input gave a result similar to direct ML. Within a region with all labels, direct ML reached a lower point error than the physics-residual model in the Lena Delta (\nrmseLeblkDz{} against \nrmseLeblkRo{} cm), but its interval against Stefan included zero there and in Alaska, and it failed under region hold-out (above). Replacing CatBoost by random forest, a multilayer perceptron, TabM or FT-Transformer did not change the sign of the physics-residual gain in the Lena Delta and Canada.

Within regional terciles of eight covariates the gain appeared in most conditions of the Lena Delta (\ngcLesig{} of \ngcLecells{} terciles with an interval excluding zero) and Canada (\ngcCasig{} of \ngcCacells), and it was largest where the Stefan model had the largest error: across terciles the correlation between the Stefan error and the gain was \ngcLecorr{} in the Lena Delta and \ngcCacorr{} in Canada (Fig.~3d). In the warmest TDD tercile of the Lena Delta the gain was \ndtddLewarm{} cm. In Alaska no tercile changed by more than \ngcAkmaxabs{} cm. The residual model therefore acts as a correction of the Stefan coefficient where soil and terrain make the coefficient deviate from its regional value, which is also where its inputs carry the most weight: terrain contributed \nshapLeResTerrain\% of the mean absolute SHAP value of the residual model in the Lena Delta and soil \nshapCaResSoil\% in Canada (Supplementary Fig.~S3).

\subsection{Few labels and new regions}
The label ladders show how the methods behave when labels are scarce (Fig.~4a). With 25 training labels drawn from the 2010--2024 cell-years of a region, the physics-residual model was at least as accurate as a Stefan model recalibrated on the same labels (Alaska \nllAkRotwfive, Lena Delta \nllLeRotwfive, Canada \nllCaRotwfive{} cm), and with 100 labels its gain was resolved in all three regions (\nllAkRohundred, \nllLeRohundred{} and \nllCaRohundred{} cm). Direct ML with 25 labels had a higher point error than Stefan in Alaska (\nllAkDztwfive{} cm) and Canada (\nllCaDztwfive{} cm).

In the transfer experiment a region was predicted by models trained on the two other regions plus $n$ labels from the target region (Fig.~4b). Without target labels, direct ML gave \ntlrmseLeDszero{} cm in the Lena Delta, against \ntlrmseLePzzero{} cm for Stefan with the source coefficient. With 25 target labels the error fell to \ntlrmseLeRstwfive{} cm for the physics-residual model and \ntlrmseLeDstwfive{} cm for direct ML. In Alaska the physics-residual model gave the lowest error at every $n$ (\ntlrmseAkRszero{} cm without target labels). In Canada direct ML trained on the source regions was the most accurate (\ntlrmseCaDszero{} cm), and recalibrating the Stefan coefficient on 25 Canadian labels increased the error from \ntlrmseCaPzzero{} to \ntlrmseCaPntwfive{} cm, and the physics-residual model, which uses this anchor, followed it (\ntlrmseCaRszero{} to \ntlrmseCaRstwfive{} cm), reflecting the wide range of thaw depths among the Canadian sites.

The same procedure was applied to eight external CALM regions (\nextcells{} cells, \nyrsext{} observations) that were not used for model development (Fig.~4c and Supplementary Fig.~S4). Without local labels the physics-residual model trained on the three study regions gave \nftAllRszero{} cm on the \nftrows{} cell-years of seven regions that also have values of Wei et al., Liu et al. and ESA CCI v5. Five labels from the same external region reduced the error to \nftAllRsfive{} cm and 25 labels to \nftAllRstwfive{} cm, whereas direct ML reached \nftAllDstwfive{} cm with 25 labels. With five labels the physics-residual model was below the ESA CCI v5 product, which was not trained on ALT observations (\nftAllCcifivezero{} cm; difference \nftdRsCcifivefive{} cm). Wei et al. and Liu et al. reached \nftAllWeizero{} and \nftAllLiuzero{} cm at these sites, which are part of their training data. Most of the gain came from recalibrating the Stefan coefficient; the residual model added no further reduction at these sites (\nftdRsPnfive{} cm with five labels), where the thaw depths of West and Central Siberia and Svalbard exceed the range of the study regions.

\subsection{Comparison with published products in the same cells and years}
Products and models were compared on the cell-years at which Wei et al., Liu et al. and ESA CCI v4 all provide a value of the observation year, which restricts the comparison to \nyrscc{} (Fig.~5a). In the Lena Delta the physics-residual model gave \nccLeRoall{} cm, lower than ESA CCI v5 (\nccLeCcifiveall{} cm), Aalto et al. (\nccLeAaltoall{} cm), ESA CCI v4 (\nccLeCcifourall{} cm), Liu et al. (\nccLeLiuall{} cm) and Wei et al. (\nccLeWeiall{} cm). Away from product training sites (more than 25 km), the paired differences were \ndprodLeweiRomask{} cm against Wei et al. and \ndprodLeccifiveRomask{} cm against ESA CCI v5. In Alaska (\nccAkRoall{} against \nccAkWeiall{} and \nccAkLiuall{} cm) and Canada (\nccCaRoall{} against \nccCaWeiall{} and \nccCaLiuall{} cm) the model and the two ML products were on par.

The difference between regions follows the training data of the products. The Alaska label cells lie at a median distance of \ndistweiAkmed{} km from the training sites of Wei et al., and Liu et al. were trained on an earlier version of the same Alaska and Canada field data. The Lena Delta cells lie at a median distance of \ndistweiLemed{} km (Supplementary Fig.~S5), and both products overestimate the thin active layers of the northern delta (mean bias \nbiasLeNorthWei{} and \nbiasLeNorthLiu{} cm north of 72$^\circ$N). A model trained without any Lena Delta label was already more accurate there than Wei et al. on the same cell-years (\ntlvprmseLeRsWeizero{} against \ntlvprmseLeWeizero{} cm; \ntlvpLeRsWeizero{} cm).

Year by year, the physics-residual model had a lower error than Wei et al. and than Stefan in all \nywLen{} Lena Delta years with at least 30 cell-years (\nyrsywLe), in \nywAkRoWeiwin{} of \nywAkRoWein{} Alaska years against Wei et al. and \nywAkRoPowin{} of \nywAkRoPon{} against Stefan (\nyrsywAk), and in \nywCaRoPowin{} of \nywCaRoPon{} Canadian years against Stefan (Fig.~5b). At Alaska cells observed in several years, the correlation between observed and predicted year-to-year changes was \nrdeltaAkRo{} for the physics-residual model and \nrdeltaAkWei{} and \nrdeltaAkLiu{} for Wei et al. and Liu et al.; in the Lena Delta the products followed the year-to-year changes more closely (\nrdeltaLeWei{} and \nrdeltaLeLiu{} against \nrdeltaLeRo).

The residual model can also be placed on top of a product. With the Lena Delta labels, a residual learner added to the linearly recalibrated product of Wei et al. changed its error by \npadLeRpPcal{} cm, and the Stefan-anchored model remained more accurate than this product-anchored model (\npadLeRoRp{} cm; Supplementary Fig.~S6). In Alaska and Canada the two anchors were not distinguishable.

\subsection{Year-labelled 1 km maps}
The physics-residual models fitted on all cell-years of Alaska and of the Lena Delta were applied to 1 km grids with the ERA5-Land climate of each year from 2011 to 2024 (Fig.~6). The grids cover cells with an ESA CCI permafrost probability of at least 50\% (\nncellmapAk{} cells in Alaska and \nncellmapLe{} in the Lena Delta). The regional mean anomaly relative to the 2011--2024 mean ranged from \nanminAk{} cm in \nanminyrAk{} to \nanmaxAk{} cm in \nanmaxyrAk{} in Alaska and from \nanminLe{} cm in \nanminyrLe{} to \nanmaxLe{} cm in \nanmaxyrLe{} in the Lena Delta, with different spatial patterns in each year (Fig.~6b). In the Lena Delta the mean anomaly of the observed cells had the same sign as the predicted regional anomaly in \nanobssameLe{} of \nanobsnLe{} years with at least 50 cell-years (\nanobsyrsLe; $r$ = \nanobsrLe). In Alaska, where the observed cells change from campaign to campaign, the agreement was weaker ($r$ = \nanobsrAk) and the predicted interannual range (\nanpredrangeAk{} cm) was smaller than that of the observed cells (\nanobsrangeAk{} cm) (Fig.~6c). \naoaAk\% of the Alaska cells and \naoaLe\% of the Lena Delta cells are inside the area of applicability, and the half-width of the 90\% prediction interval is \nhwAk{} and \nhwLe{} cm. Cells outside the area of applicability, mostly mountains and the southern margin of the permafrost zone, are marked in the maps (Supplementary Fig.~S7). Supplementary Fig.~S8 compares our maps with those of Wei et al. and Liu et al. for Alaska in 2014 and the Lena Delta in 2016, and Supplementary Fig.~S9 shows the correction that the residual model adds to the Stefan model.
